\section{Limitations}

ECHOREPO was developed within a specific soil citizen-science programme, and several components remain closely linked to that scientific and organisational context. In particular, the current canonical data model contains soil-specific concepts and the ingestion layer reflects the characteristics of the EchoSoil acquisition environment. Reuse in another domain would therefore require adaptation of the source-specific ingestion logic and at least part of the canonical schema.

The study also evaluates a single production deployment rather than multiple independent implementations of the architecture. Consequently, the results support claims of architectural reusability and operational feasibility, but do not yet demonstrate that ECHOREPO can be transferred to another citizen-science domain without substantial modification. A deployment in a second scientific context would provide stronger evidence regarding the effort required for reuse.

The quantitative evaluation is based primarily on operational data generated during the ECHO deployment rather than on a controlled comparison with alternative repository architectures. The observed data-quality outcomes therefore describe the behaviour of the ECHOREPO workflow in its actual use context, but should not be interpreted as evidence that the same error rates would occur in other citizen-science projects. Similarly, recovery rates depend on the provenance information available in this particular workflow and may differ substantially in systems that retain less contextual information.

The geographical quality-assurance analysis currently provides the most mature quantitative example of repository-level validation. Other data-quality dimensions, including semantic inconsistencies, incomplete observations and heterogeneous categorical values, are handled operationally but are not yet characterised with the same level of quantitative detail. The evaluation should therefore not be interpreted as a comprehensive assessment of all possible data-quality problems in citizen-generated environmental data.

The privacy approach also has limitations. Spatial obfuscation reduces the precision of publicly released coordinates, but it does not constitute a formal guarantee of anonymity. Residual disclosure risk depends on the spatial context, the availability of auxiliary information and the relationship between sampling locations and participants. The current mechanism should therefore be understood as a pragmatic privacy-preserving publication measure rather than a complete privacy model.

Likewise, the FAIR-oriented publication workflow demonstrates several concrete mechanisms supporting findability, accessibility, interoperability and reuse, including persistent identifiers, machine-readable resources, structured metadata and external discovery. However, the study does not claim complete FAIR compliance. In particular, semantic interoperability remains dependent on the availability and quality of controlled vocabularies, metadata mappings and externally recognised identifiers.

ECHOREPO is also an evolving research infrastructure. Automated test coverage, external deployment documentation and the
separation between generic infrastructure and project-specific
configuration are being improved, and some parts of the current implementation reflect historical design decisions made while the system was already supporting an active project. This may limit immediate software reuse even where the underlying architectural pattern is applicable.

Finally, several quantitative measures reported in this study are snapshots of a continuously changing repository. Sample counts, storage volume, enrichment coverage and published dataset versions may therefore evolve after manuscript preparation. For this reason, the evaluation should report the date or version associated with each quantitative snapshot and, where possible, link the results to archived research-data releases and reproducible analysis scripts.

These limitations define the scope of the conclusions. The present study provides evidence that the ECHOREPO architecture can support heterogeneous, quality-controlled and FAIR-oriented management of citizen-science soil data at operational scale. It does not establish that the implementation is universally applicable without adaptation, nor that the specific policies and validation rules used in ECHO are optimal for other projects. The broader contribution lies in the design principles and operational lessons that can be evaluated and adapted in other citizen-science research infrastructures.